% Canonical LaTeX source; imported and revised from the legacy Markdown.
\chapter{Logica, insiemi, numeri e strumenti discreti}
\label{chap:1}

\begin{obiettivocapitolo}{usare i prerequisiti come strumenti di calcolo e controllo}
\prioritaalta. Le competenze di questo capitolo intervengono in quasi ogni esercizio: segni, frazioni, potenze, valore assoluto, sommatorie e stime.
\end{obiettivocapitolo}
\begin{sceltametodo}{riconoscere l'operazione necessaria}
\begin{tabularx}{\textwidth}{@{}p{.30\textwidth}X@{}}
Frazioni o rapporti & porta a denominatore comune; semplifica solo fattori, mai addendi.\\
Percentuali successive & moltiplica i fattori $1+p_j/100$; non sommare automaticamente le percentuali.\\
Potenze e radicali & verifica parità dell'indice, segno della base e possibili rami.\\
Sommatorie & separa linearmente, cambia indice con coerenza, riconosci progressioni.\\
Valore assoluto & traduci in distanza o separa i casi in base al segno dell'argomento.\\
Stima & usa disuguaglianza triangolare, Cauchy--Schwarz, AM--GM o confronto con una quantità più semplice.
\end{tabularx}
\end{sceltametodo}
\begin{errorefrequente}{semplificazioni illegittime}
È lecito cancellare un fattore comune non nullo, non un termine di una somma: $\frac{x(x+1)}x=x+1$ per $x\ne0$, mentre $\frac{x+1}{x}$ non si semplifica cancellando $x$.
\end{errorefrequente}
\begin{controllofinale}{prerequisiti}
Controlla segno, denominatori, ordine di grandezza, unità di misura e coerenza fra valore esatto e approssimato.
\end{controllofinale}

\section{Logica matematica minima}\label{logica-matematica-minima}

Nelle formule seguenti \(P\) e \(Q\) sono proposizioni, mentre \(P(x)\) è una proprietà che può essere vera o falsa a seconda di \(x\). I simboli \(\mathrm V\) e \(\mathrm F\) indicano rispettivamente «vero» e «falso».

\[
\neg(\forall x\ P(x))\Longleftrightarrow\exists x:\neg P(x),
\qquad
\neg(\exists x\ P(x))\Longleftrightarrow\forall x:\neg P(x).
\]

\[
P\Rightarrow Q
\Longleftrightarrow
\neg P\lor Q
\Longleftrightarrow
\neg Q\Rightarrow\neg P.
\]

L'ultima equivalenza è la \textbf{contrapposizione}: «se \(P\) implica \(Q\), allora non \(Q\) implica non \(P\)». Non va confusa con il converso \(Q\Rightarrow P\).

\[
P\Longleftrightarrow Q
\Longleftrightarrow
(P\Rightarrow Q)\land(Q\Rightarrow P).
\]

\[
\begin{array}{c|c|c|c}
P&Q&P\Rightarrow Q&P\Longleftrightarrow Q\\
\hline
\mathrm V&\mathrm V&\mathrm V&\mathrm V\\
\mathrm V&\mathrm F&\mathrm F&\mathrm F\\
\mathrm F&\mathrm V&\mathrm V&\mathrm F\\
\mathrm F&\mathrm F&\mathrm V&\mathrm V
\end{array}
\]

Approfondimento: \href{https://ingegnerismo.it/matematica/logica-matematica/}{logica matematica}, \href{https://ingegnerismo.it/matematica/quantificatori/}{quantificatori}.

\section{Insiemi, intervalli e intorni}\label{insiemi-intervalli-e-intorni}

Nella notazione \(\{x:P(x)\}\) i due punti si leggono «tale che». In questa sottosezione \(U\) è l'insieme universo rispetto al quale si prende il complementare.

\[
A\cup B=\{x:x\in A\lor x\in B\},
\qquad
A\cap B=\{x:x\in A\land x\in B\},
\]

\[
A\setminus B=\{x\in A:x\notin B\},
\qquad
A^c=U\setminus A.
\]

\begin{longtable}[]{@{}ll@{}}
\toprule\noalign{}
Intervallo & Condizione su \(x\) \\
\midrule\noalign{}
\endhead
\bottomrule\noalign{}
\endlastfoot
\((a,b)\) & \(a<x<b\) \\
\([a,b]\) & \(a\le x\le b\) \\
\((a,b]\) & \(a<x\le b\) \\
\([a,b)\) & \(a\le x<b\) \\
\((a,+\infty)\) & \(x>a\) \\
\([a,+\infty)\) & \(x\ge a\) \\
\((-\infty,b)\) & \(x<b\) \\
\((-\infty,b]\) & \(x\le b\) \\
\end{longtable}

\[
I(x_0,r)=(x_0-r,x_0+r),
\qquad
I^*(x_0,r)=I(x_0,r)\setminus\{x_0\},
\qquad r>0.
\]

\[
I_D(x_0,r)=D\cap I(x_0,r),
\qquad
I_D^*(x_0,r)=D\cap I^*(x_0,r).
\]

\[
x_0\in\operatorname{int}A
\Longleftrightarrow
\exists r>0:\ I(x_0,r)\subseteq A.
\]

\[
x_0\in A'
\Longleftrightarrow
\forall r>0:\ I^*(x_0,r)\cap A\ne\varnothing.
\]

\[
x_0\notin A'
\Longleftrightarrow
\exists r>0:\ I^*(x_0,r)\cap A=\varnothing.
\]

\[
a\text{ isolato in }A
\Longleftrightarrow
a\in A\ \land\ \exists r>0:\ I(a,r)\cap A=\{a\}.
\]

\[
\overline A=A\cup A'.
\]

Qui \(A'\) è l'insieme dei punti di accumulazione, \(\overline A\) è la chiusura e \(\operatorname{int}A\) è l'interno. Gli intorni relativi \(I_D\) limitano la ricerca ai soli punti appartenenti al dominio \(D\).

Approfondimento: \href{https://ingegnerismo.it/matematica/insieme/}{insieme}, \href{https://ingegnerismo.it/matematica/operazioni-tra-insiemi/}{operazioni tra insiemi}, \href{https://ingegnerismo.it/matematica/intervallo-reale/}{intervallo reale}, \href{https://ingegnerismo.it/matematica/intorni-e-topologia-della-retta/}{intorni e topologia della retta}, \href{https://ingegnerismo.it/matematica/punto-di-accumulazione/}{punto di accumulazione}.

Esercizi svolti: \href{https://ingegnerismo.it/matematica/intorni-accumulazione-esercizi/}{intorni, punti di accumulazione e punti isolati}.

\section{Topologia essenziale della retta reale}\label{topologia-essenziale-della-retta-reale}

Per \(A\subseteq\mathbb R\):

\[
A^\circ=\operatorname{int}A,
\qquad
\partial A=\overline A\setminus A^\circ
=\overline A\cap\overline{A^c}.
\]

\[
\begin{aligned}
A\text{ aperto}
&\Longleftrightarrow A=A^\circ,\\
A\text{ chiuso}
&\Longleftrightarrow A'\subseteq A
\Longleftrightarrow A=\overline A.
\end{aligned}
\]

\[
A\text{ limitato}
\Longleftrightarrow
\exists M>0:\ |x|\le M\quad\forall x\in A.
\]

In \(\mathbb R\) vale il teorema di Heine--Borel:

\[
K\text{ compatto}
\Longleftrightarrow
K\text{ chiuso e limitato}.
\]

Caratterizzazione sequenziale della compattezza:

\[
K\text{ compatto}
\Longleftrightarrow
\text{ogni successione in }K\text{ ammette una sottosuccessione convergente a un punto di }K.
\]

Un sottoinsieme non vuoto di \(\mathbb R\) è connesso se e solo se è un intervallo.

Cardinalità essenziale:

\[
\mathbb N_0,\ \mathbb Z,\ \mathbb Q\text{ sono numerabili},
\qquad
\mathbb R\text{ non è numerabile}.
\]

Per l'insieme delle parti

\[
\mathcal P(A)=\{B:B\subseteq A\},
\]

il teorema di Cantor afferma

\[
\operatorname{card}\mathcal P(A)>\operatorname{card}A.
\]

\section{Insiemi numerici, ordine e completezza}\label{insiemi-numerici-ordine-e-completezza}

\[
\mathbb N_0\subset\mathbb Z\subset\mathbb Q\subset\mathbb R\subset\mathbb C,
\qquad
\mathbb N_0=\{0,1,2,\ldots\},
\qquad
\mathbb N=\mathbb N_+=\{1,2,3,\ldots\}.
\]

\[
M=\max A
\Longleftrightarrow
M\in A\ \land\ \forall x\in A:\ x\le M,
\]

\[
m=\min A
\Longleftrightarrow
m\in A\ \land\ \forall x\in A:\ m\le x.
\]

Per \(A\ne\varnothing\) superiormente limitato:

\[
s=\sup A
\Longleftrightarrow
\begin{cases}
\forall x\in A,\ x\le s,\\
\forall\varepsilon>0,\ \exists x\in A:\ s-\varepsilon<x\le s.
\end{cases}
\]

Per \(A\ne\varnothing\) inferiormente limitato:

\[
i=\inf A
\Longleftrightarrow
\begin{cases}
\forall x\in A,\ i\le x,\\
\forall\varepsilon>0,\ \exists x\in A:\ i\le x<i+\varepsilon.
\end{cases}
\]

\[
A\ne\varnothing,\ A\subseteq\mathbb R,\ A\text{ superiormente limitato}
\Longrightarrow
\sup A\in\mathbb R.
\]

\[
A\ne\varnothing,\ A\subseteq\mathbb R,\ A\text{ inferiormente limitato}
\Longrightarrow
\inf A\in\mathbb R.
\]

Per \(A\ne\varnothing\) superiormente limitato:

\[
\max A\text{ esiste}
\Longleftrightarrow
\sup A\in A.
\]

Per \(A\ne\varnothing\) inferiormente limitato:

\[
\min A\text{ esiste}
\Longleftrightarrow
\inf A\in A.
\]

Il massimo e il minimo, quando esistono, appartengono all'insieme; il supremo e l'infimo possono non appartenervi.

\[
\forall x\in\mathbb R,\ \exists n\in\mathbb N_+:\ n>x.
\]

\[
x<y
\Longrightarrow
\exists q\in\mathbb Q:\ x<q<y,
\qquad
\exists r\in\mathbb R\setminus\mathbb Q:\ x<r<y.
\]

Approfondimento: \href{https://ingegnerismo.it/matematica/insiemi-numerici/}{insiemi numerici}, \href{https://ingegnerismo.it/matematica/completezza-dei-reali/}{completezza dei reali}, \href{https://ingegnerismo.it/matematica/estremo-superiore/}{estremo superiore}.

\section{Aritmetica, divisibilità e segni}\label{aritmetica-divisibilituxe0-e-segni}

La tabella seguente vale sia per la moltiplicazione sia per la divisione, purché il divisore sia non nullo.

\[
\begin{array}{c|cc}
\cdot\text{ o }:&+&-\\
\hline
+&+&-\\
-&-&+
\end{array}
\]

\[
\frac{a+b}{c-d}=(a+b):(c-d),
\qquad c-d\ne0.
\]

\[
a\mid b
\Longleftrightarrow
\exists k\in\mathbb Z:\ b=ak.
\]

Per \(a,b\in\mathbb Z\) con \(a\ne0\) esistono unici \(q,r\in\mathbb Z\) tali che

\[
b=aq+r,
\qquad
0\le r<|a|.
\]

Se \(a,b\in\mathbb N_+\) sono scritti mediante la fattorizzazione prima

\[
a=\prod_p p^{\alpha_p},
\qquad
b=\prod_p p^{\beta_p},
\]

allora

\[
\gcd(a,b)=\prod_p p^{\min(\alpha_p,\beta_p)},
\qquad
\operatorname{lcm}(a,b)=\prod_p p^{\max(\alpha_p,\beta_p)}.
\]

\[
\gcd(a,b)\operatorname{lcm}(a,b)=ab,
\qquad a,b>0.
\]

Approfondimento: \href{https://ingegnerismo.it/matematica/aritmetica-interi-frazioni-decimali/}{aritmetica di interi, frazioni e decimali}.

\section{Frazioni, decimali e notazione scientifica}\label{frazioni-decimali-e-notazione-scientifica}

\[
\frac ab=\frac cd
\Longleftrightarrow
ad=bc,
\qquad b,d\ne0.
\]

\[
\frac{-a}{b}=\frac{a}{-b}=-\frac ab,
\qquad
\frac{-a}{-b}=\frac ab.
\]

\[
\frac ab\pm\frac cd=\frac{ad\pm bc}{bd},
\qquad
\frac ab\frac cd=\frac{ac}{bd},
\]

\[
\frac{a/b}{c/d}=\frac{ad}{bc},
\qquad
b,c,d\ne0.
\]

\[
\frac{ac}{bc}=\frac ab\quad(bc\ne0).
\]

\[
\begin{array}{c|c}
\text{forma}&\text{valore}\\
\hline
0/b\ (b\ne0)&0\\
a/0&\text{non definita}\\
0/0&\text{non definita}
\end{array}
\]

Per \(a/b\) ridotta ai minimi termini, con \(b>0\):

\[
\frac ab\text{ ha sviluppo decimale finito}
\Longleftrightarrow
b=2^m5^n,
\qquad m,n\in\mathbb N_0.
\]

\[
x=m\,10^k,
\qquad
1\le|m|<10,
\qquad
k\in\mathbb Z,
\qquad x\ne0.
\]

\[
(m_1 10^{k_1})(m_2 10^{k_2})=(m_1m_2)10^{k_1+k_2},
\]

\[
\frac{m_1 10^{k_1}}{m_2 10^{k_2}}
=\frac{m_1}{m_2}10^{k_1-k_2},
\qquad m_2\ne0.
\]

Approfondimento: \href{https://ingegnerismo.it/matematica/frazioni-numeriche/}{frazioni numeriche}, \href{https://ingegnerismo.it/matematica/notazione-scientifica-approssimazioni/}{notazione scientifica e approssimazioni}.

\section{Rapporti, proporzioni e percentuali}\label{rapporti-proporzioni-e-percentuali}

\textbf{Proporzione.} Nella scrittura \(a:b=c:d\), i termini \(b\) e \(d\) sono i conseguenti e devono essere non nulli:

\[
a:b=c:d
\Longleftrightarrow
\frac ab=\frac cd
\Longleftrightarrow
ad=bc,
\qquad b,d\ne0.
\]

Forme risolte più usate:

\begin{longtable}[]{@{}lll@{}}
\toprule\noalign{}
Incognita & Formula & Condizioni operative \\
\midrule\noalign{}
\endhead
\bottomrule\noalign{}
\endlastfoot
\(a\) & \(a=\dfrac{bc}{d}\) & \(d\ne0\) \\
\(c\) & \(c=\dfrac{ad}{b}\) & \(b\ne0\) \\
\(b\) & \(b=\dfrac{ad}{c}\) & \(c\ne0\), \(d\ne0\) e risultato \(b\ne0\) \\
\(d\) & \(d=\dfrac{bc}{a}\) & \(a\ne0\), \(b\ne0\) e risultato \(d\ne0\) \\
\end{longtable}

Se \(a=c=0\), ogni coppia \(b,d\ne0\) soddisfa la proporzione e né \(b\) né \(d\) è determinato univocamente. Se uno solo tra \(a\) e \(c\) è nullo, la proporzione non può essere verificata con \(b,d\ne0\).

\textbf{Proporzionalità diretta.} \(x\) è la variabile indipendente, \(y\) quella dipendente e \(k\) la costante di proporzionalità:

\[
y=kx,
\qquad
x=\frac yk\quad(k\ne0),
\qquad
k=\frac yx\quad(x\ne0).
\]

Se \(k=0\), allora \(y=0\) per ogni \(x\) e la relazione non permette di ricavare un unico \(x\) da \(y\).

\textbf{Proporzionalità inversa:}

\[
y=\frac kx,
\qquad x\ne0,
\]

\[
k=xy,
\qquad
x=\frac ky\quad(y\ne0).
\]

Se \(k=0\), si ha \(y=0\) per ogni \(x\ne0\); anche in questo caso non esiste una formula inversa univoca per \(x\).

\textbf{Percentuali.} \(Q_i\) è il valore iniziale, \(Q_f\) il valore finale e \(p\) la variazione percentuale \textbf{con segno}: \(p>0\) indica un aumento, \(p<0\) una riduzione.

\[
p\%=\frac p{100}.
\]

Per calcolare il \(p\%\) di una quantità \(Q\), posto \(V\) il valore percentuale:

\[
V=Q\frac p{100},
\qquad
Q=\frac{100V}{p}\quad(p\ne0),
\qquad
p=\frac{100V}{Q}\quad(Q\ne0).
\]

Se \(p=0\), allora \(V=0\) per ogni \(Q\) e \(Q\) non è ricavabile in modo univoco; se \(Q=0\), allora \(V=0\) per ogni \(p\) e \(p\) non è ricavabile in modo univoco.

Formula diretta e forme inverse per una variazione percentuale:

\[
Q_f=Q_i\left(1+\frac p{100}\right),
\]

\[
Q_i=\frac{Q_f}{1+p/100},
\qquad p\ne-100,
\]

\[
p=100\left(\frac{Q_f}{Q_i}-1\right),
\qquad Q_i\ne0.
\]

Se si preferisce indicare una riduzione mediante una percentuale positiva \(r\ge0\):

\[
Q_f=Q_i\left(1-\frac r{100}\right),
\qquad
Q_i=\frac{Q_f}{1-r/100}\quad(r\ne100),
\]

\[
r=100\left(1-\frac{Q_f}{Q_i}\right),
\qquad Q_i\ne0.
\]

Variazioni successive \(p_1,\ldots,p_m\):

\[
Q_f=Q_i\prod_{j=1}^m\left(1+\frac{p_j}{100}\right),
\]

\[
p_{\mathrm{tot}}
=100\left[
\prod_{j=1}^m\left(1+\frac{p_j}{100}\right)-1
\right].
\]

La percentuale che riporta \(Q_f\) al valore iniziale non è, in generale, \(-p\). Se \(p\) è la variazione da \(Q_i\) a \(Q_f\), la variazione inversa è

\[
p_{\mathrm{inv}}=-\frac{100p}{100+p},
\qquad p\ne-100.
\]

In particolare, applicare prima \(+p\%\) e poi \(-p\%\) produce

\[
\left(1+\frac p{100}\right)\left(1-\frac p{100}\right)
=1-\left(\frac p{100}\right)^2,
\]

quindi non si torna al valore iniziale, salvo \(p=0\).

Approfondimento: \href{https://ingegnerismo.it/matematica/rapporti-proporzioni-percentuali/}{rapporti, proporzioni e percentuali}, \href{https://ingegnerismo.it/matematica/proporzionalita-diretta-e-inversa/}{proporzionalità diretta e inversa}, \href{https://ingegnerismo.it/matematica/variazioni-percentuali/}{variazioni percentuali}.

\section{Potenze e radicali reali}\label{potenze-e-radicali-reali}

Per \(n\in\mathbb N_+\):

\[
a^n=\underbrace{a\cdot\ldots\cdot a}_{n\text{ fattori}},
\qquad
a^0=1\ (a\ne0),
\qquad
a^{-n}=\frac1{a^n}\ (a\ne0).
\]

\[
\begin{array}{c|c}
\text{forma}&\text{valore nel campo reale}\\
\hline
0^n,\ n>0&0\\
0^0&\text{non definita}\\
0^{-n},\ n>0&\text{non definita}
\end{array}
\]

Per \(m,n\in\mathbb Z\), quando tutte le potenze sono definite (in particolare \(a\ne0\) se compare un esponente negativo):

\[
a^ma^n=a^{m+n},
\qquad
\frac{a^m}{a^n}=a^{m-n}\quad(a\ne0),
\]

\[
(a^m)^n=a^{mn},
\qquad
(ab)^n=a^nb^n,
\qquad
\left(\frac ab\right)^n=\frac{a^n}{b^n}\ (b\ne0).
\]

La radice \(\sqrt[n]{a}\) indica il \textbf{ramo reale principale}. Per \(n\in\mathbb N_+\):

\[
\sqrt[n]{a}=r
\Longleftrightarrow
\begin{cases}
r^n=a\ \land\ r\ge0,&n\text{ pari e }a\ge0,\\
r^n=a,&n\text{ dispari e }a\in\mathbb R.
\end{cases}
\]

L'equazione \(x^n=a\) è invece una relazione da risolvere e può avere più rami:

\begin{longtable}[]{@{}lll@{}}
\toprule\noalign{}
Parità di \(n\) & Condizione su \(a\) & Soluzioni reali di \(x^n=a\) \\
\midrule\noalign{}
\endhead
\bottomrule\noalign{}
\endlastfoot
\(n\) pari & \(a>0\) & \(x=\pm\sqrt[n]{a}\) \\
\(n\) pari & \(a=0\) & \(x=0\) \\
\(n\) pari & \(a<0\) & nessuna \\
\(n\) dispari & \(a\in\mathbb R\) & \(x=\sqrt[n]{a}\) \\
\end{longtable}

\[
\sqrt{a^2}=|a|,
\qquad
(\sqrt[n]{a})^n=a,
\]

\[
\sqrt[n]{a^n}=
\begin{cases}
|a|,&n\text{ pari},\\
a,&n\text{ dispari}.
\end{cases}
\]

Per \(m\in\mathbb Z\), \(n\in\mathbb N_+\) e \(m/n\) ridotto ai minimi termini:

\[
a^{m/n}=(\sqrt[n]{a})^m.
\]

\[
\begin{array}{c|c}
n\text{ pari}&a\ge0;\ a>0\text{ se }m<0\\
n\text{ dispari}&a\in\mathbb R;\ a\ne0\text{ se }m<0
\end{array}
\]

\[
a^x=e^{x\ln a},
\qquad a>0,
\qquad x\in\mathbb R.
\]

\[
\begin{array}{c|c|c}
&\sqrt[n]{ab}=\sqrt[n]a\sqrt[n]b
&\sqrt[n]{a/b}=\sqrt[n]a/\sqrt[n]b\\
\hline
n\text{ dispari}&a,b\in\mathbb R&a\in\mathbb R,\ b\ne0\\
n\text{ pari}&a,b\ge0&a\ge0,\ b>0
\end{array}
\]

\[
\frac1{\sqrt a}=\frac{\sqrt a}{a},
\qquad a>0,
\]

\[
\frac1{\sqrt a+\sqrt b}
=\frac{\sqrt a-\sqrt b}{a-b},
\qquad a,b\ge0,
\qquad a\ne b.
\]

\[
x^2=a
\Longleftrightarrow
x=\pm\sqrt a,
\qquad a\ge0,
\]

\[
\sqrt{x^2}=\lvert x\rvert.
\]

Approfondimento: \href{https://ingegnerismo.it/matematica/potenze-e-radicali/}{potenze e radicali}.

\section{Principio di induzione e disuguaglianza di Bernoulli}\label{principio-di-induzione-e-disuguaglianza-di-bernoulli}

\[
\left.
\begin{array}{l}
P(n_0)\\
\forall n\ge n_0:\ P(n)\Rightarrow P(n+1)
\end{array}
\right\}
\Longrightarrow
\forall n\ge n_0:\ P(n).
\]

\[
(1+x)^n\ge1+nx,
\qquad
n\in\mathbb N_+,
\qquad
x\ge-1.
\]

\[
(1+x)^n=1+nx
\Longleftrightarrow
n=1\ \lor\ x=0.
\]

Approfondimento: \href{https://ingegnerismo.it/matematica/principio-di-induzione/}{principio di induzione}, \href{https://ingegnerismo.it/matematica/disuguaglianza-di-bernoulli/}{disuguaglianza di Bernoulli}.

\section{Fattoriale e coefficienti binomiali}\label{fattoriale-e-coefficienti-binomiali}

\[
n!=\prod_{k=1}^n k,
\qquad
0!=1,
\qquad
n\in\mathbb N_0.
\]

\[
\binom nk=\frac{n!}{k!(n-k)!},
\qquad
n\in\mathbb N_0,
\qquad
0\le k\le n.
\]

\[
\binom nk=\binom n{n-k},
\qquad 0\le k\le n,
\]

\[
\binom nk+\binom n{k+1}=\binom{n+1}{k+1},
\qquad 0\le k\le n-1.
\]

\[
(a+b)^n=\sum_{k=0}^n\binom nk a^{n-k}b^k.
\]

Approfondimento: \href{https://ingegnerismo.it/matematica/fattoriale/}{fattoriale}, \href{https://ingegnerismo.it/matematica/coefficiente-binomiale/}{coefficiente binomiale}.

\section{Sommatorie, progressioni e somme finite}\label{sommatorie-progressioni-e-somme-finite}

Nella sommatoria \(\sum_{k=m}^{n}a_k\), \(k\) è un \textbf{indice muto}: può essere rinominato senza cambiare il valore, purché siano modificati coerentemente estremi e termine generale.

Per \(m,n\in\mathbb Z\) con \(m\le n\), linearità e cambio di indice:

\[
\sum_{k=m}^{n}(\alpha a_k+\beta b_k)
=
\alpha\sum_{k=m}^{n}a_k
+
\beta\sum_{k=m}^{n}b_k,
\]

\[
\sum_{k=m}^{n}a_k
=
\sum_{j=0}^{n-m}a_{m+j}.
\]

Somme fondamentali, per \(n\in\mathbb N_+\):

\[
\sum_{k=1}^{n}1=n,
\qquad
\sum_{k=1}^{n}k=\frac{n(n+1)}2,
\]

\[
\sum_{k=1}^{n}k^2=\frac{n(n+1)(2n+1)}6,
\qquad
\sum_{k=1}^{n}k^3=\left[\frac{n(n+1)}2\right]^2.
\]

\textbf{Progressione aritmetica.} \(a_1\) è il primo termine, \(d\) la differenza comune, \(a_n\) il termine di indice \(n\) e \(S_n\) la somma dei primi \(n\) termini:

\[
a_n=a_1+(n-1)d,
\]

\[
S_n=\sum_{k=1}^{n}a_k
=\frac n2(a_1+a_n)
=\frac n2\bigl[2a_1+(n-1)d\bigr].
\]

Forme inverse principali, per \(n>1\):

\[
d=\frac{a_n-a_1}{n-1},
\qquad
a_1=a_n-(n-1)d,
\]

\[
a_1=\frac{S_n}{n}-\frac{n-1}{2}d,
\qquad
d=\frac{2(S_n-na_1)}{n(n-1)}.
\]

Se \(d\ne0\):

\[
n=1+\frac{a_n-a_1}{d},
\]

ma il risultato è un indice valido soltanto se appartiene a \(\mathbb N_+\).

\textbf{Progressione geometrica.} \(q\) è la ragione comune. Il primo termine è sempre \(a_1\); per \(n>1\):

\[
a_n=a_1q^{n-1}.
\]

\[
S_n=\sum_{k=1}^{n}a_k
=
\begin{cases}
\displaystyle a_1\frac{1-q^n}{1-q},&q\ne1,\\[8pt]
na_1,&q=1.
\end{cases}
\]

Forme inverse utili:

\[
a_1=a_n\quad(n=1),
\qquad
a_1=\frac{a_n}{q^{n-1}}\quad(n>1,\ q\ne0),
\]

\[
q=\frac{a_{n+1}}{a_n}\quad(a_n\ne0),
\]

\[
q^{n-1}=\frac{a_n}{a_1}\quad(a_1\ne0).
\]

L'ultima relazione può produrre più rami reali o nessun ramo, a seconda del segno di \(a_n/a_1\) e della parità di \(n-1\). Se \(q>0\), \(q\ne1\) e \(a_n/a_1>0\), si può ricavare anche l'indice:

\[
n=1+\frac{\ln(a_n/a_1)}{\ln q},
\]

accettando il risultato soltanto quando \(n\in\mathbb N_+\).

\section{Valore assoluto}\label{valore-assoluto}

\[
|x|=
\begin{cases}
x,&x\ge0,\\
-x,&x<0.
\end{cases}
\]

\[
|x|\ge0,
\qquad
|x|=0\Longleftrightarrow x=0,
\qquad
|-x|=|x|.
\]

\[
|xy|=|x||y|,
\qquad
\left|\frac xy\right|=\frac{|x|}{|y|}\quad(y\ne0).
\]

\[
|x+y|\le|x|+|y|,
\qquad
\bigl||x|-|y|\bigr|\le|x-y|.
\]

Per \(r\ge0\):

\[
|x-a|<r\Longleftrightarrow a-r<x<a+r,
\]

\[
|x-a|\le r\Longleftrightarrow a-r\le x\le a+r,
\]

\[
|x-a|>r\Longleftrightarrow x<a-r\ \lor\ x>a+r,
\]

\[
|x-a|\ge r\Longleftrightarrow x\le a-r\ \lor\ x\ge a+r.
\]

Approfondimento: \href{https://ingegnerismo.it/matematica/valore-assoluto/}{valore assoluto}, \href{https://ingegnerismo.it/matematica/disuguaglianza-triangolare/}{disuguaglianza triangolare}.

\section{Disuguaglianze fondamentali}\label{disuguaglianze-fondamentali}

Per \(a,b\in\mathbb R\):

\[
2|ab|\le a^2+b^2,
\qquad
(a-b)^2\ge0.
\]

Media aritmetica e geometrica, per \(a,b\ge0\):

\[
\sqrt{ab}\le\frac{a+b}{2},
\]

con uguaglianza se e solo se \(a=b\).

Disuguaglianza di Cauchy--Schwarz in \(\mathbb R^n\):

\[
\left(\sum_{k=1}^{n}a_kb_k\right)^2
\le
\left(\sum_{k=1}^{n}a_k^2\right)
\left(\sum_{k=1}^{n}b_k^2\right).
\]

L'uguaglianza vale se e solo se i vettori \((a_1,\ldots,a_n)\) e \((b_1,\ldots,b_n)\) sono linearmente dipendenti, compreso il caso in cui uno dei due sia nullo.
